\documentclass[11pt]{article}
\usepackage{jtsguide}
\setshorttitle{Technical Specifications}

\title{JTS Topology Suite --- Technical Specifications}
\author{M.\ Davis / J.\ Aquino (v1.4 / 1.4.1) \\ rebuilt from the 2003 text}
\date{}

\begin{document}
\jtscover{Technical Specifications}{Version 1.4 / 1.4.1, rebuilt}{Base text: 2003 \textperiodcentered\ rebuild: 16 August 2026}

\noindent
This document is the design specification for the classes, methods
and algorithms implemented in the JTS Topology Suite. The 2003
chapter plan is kept. Sections that would now lie about the spatial
model, WKT/WKB, overlay, or Hausdorff are amended against the code
on this tree.

Package names in this tree are \texttt{org.locationtech.jts}. The
2003 copies used \texttt{com.vividsolutions.jts}. LocationTech JTS
remains the upstream project; this fork is not described as if it
were already upstream.

\clearpage
\changectrl{%
  1.3 & 31 Mar 2003 & M.\ Davis & Updated to cover changes in JTS 1.3. \\
  1.4 & 17 Oct 2003 & J.\ Aquino & CoordinateSequences and the user-data field. \\
  1.4.1 & 2003 & M.\ Davis & Fixed definition of contains. \\
  1.4-rebuild & 16 Aug 2026 & This tree & LaTeX rebuild. SFS Curve /
    MultiCurve / MultiSurface are no longer abstract-only. WKB 8--12,
    OverlayNGCurve, hulls, and the two D-HF pairs are recorded. \\
}

\tableofcontents
\clearpage

\section{Overview}
The JTS Topology Suite is a Java API that implements a core set of
spatial data operations using an explicit precision model and robust
geometric algorithms. JTS is intended to be used in the development
of applications that support the validation, cleaning, integration
and querying of spatial datasets.

JTS attempts to implement the OpenGIS Simple Features Specification
(SFS) as accurately as possible. Where the SFS is unclear or omits a
specification, JTS chooses a reasonable and consistent alternative.
Differences from and elaborations of the SFS are documented here.

The 2003 spec stopped at linear SFS. This tree adds opt-in
ISO 19125-2 / SQL-MM curve types in \texttt{jts-curve}. The detailed
class hierarchy remains JavaDoc. This rebuild records the contract a
reader of the 2003 PDF would otherwise get wrong.

\section{Other resources}
\begin{itemize}
  \item OpenGIS Simple Features Specification for SQL Revision 1.1
        (SFS) --- master specification for the linear spatial model
        and predicates.
  \item OGC Simple Features Access 1.2.1 / ISO 19125-2 --- curve
        types and WKB codes 8--12.
  \item \emph{JTS Developer's Guide} and \emph{TestBuilder \&
        TestRunner User Guide} (siblings in this directory), rebuilt
        against the same tree.
\end{itemize}

\section{Design goals}
The 2003 goals are still the goals of JTS core:
\begin{itemize}
  \item Conform to the SFS as accurately as possible, consistent with
        a correct implementation.
  \item Follow Java conventions (\texttt{getX} / \texttt{setX},
        \texttt{isX}, lowercase method names).
  \item Support a user-defined precision model; algorithms are robust
        under it.
  \item Return topologically and geometrically correct results within
        that precision model wherever possible.
  \item Correctness is the highest priority; space and time
        efficiency is important but secondary.
  \item Fast enough for production. Algorithms and code should stay
        clear.
\end{itemize}

The curve module adds a further constraint that is not in the 2003
spec: a closed-form curve path that would be slower than
densify-then-core is refused, and the linear path runs instead.

\section{Terminology}
\begin{center}
\begin{tabular}{>{\raggedright\arraybackslash}p{3.4cm}p{9.8cm}}
  \toprule
  \textbf{Term} & \textbf{Definition} \\
  \midrule
  Coordinate     & A point in space exactly representable under the
                   precision model. \\
  Node           & A point where two lines intersect. Not necessarily
                   a representable coordinate. \\
  Noding         & Computing the nodes where geometries intersect. \\
  Robust computation & Numerical computation guaranteed to return
                   the correct answer for all inputs. \\
  SFS            & OGC Simple Features Specification (linear). \\
  SQL-MM / SFA 1.2.1 & Extended types: CircularString, CompoundCurve,
                   CurvePolygon, MultiCurve, MultiSurface. \\
  Closed form    & A cheap shape check that answers exactly. On this
                   tree: certified discs, the single-arc hull, the
                   compound-curve hull of circular-plus-straight
                   members, and the two D-HF pairs. \\
  Chord / densify & The fallback: replace arcs with LineString
                   segments. Not a laser. \\
  Vertex         & A corner point stored explicitly. \\
  \bottomrule
\end{tabular}
\end{center}

\section{Notation}
Items that adhere to the SFS are marked (SFS). Items that elaborate
or differ are marked (JTS). Items that exist only on this tree's
curve module are marked (jts-curve).

\section{Java implementation}
Where SFS coding style differs from Java, JTS follows Java:
\texttt{boolean} instead of Integer-as-boolean; method names start
with a lowercase letter; get/set prefixes for JavaBeans. The 2003
package root \texttt{com.vividsolutions.jts} is
\texttt{org.locationtech.jts} on this tree. That is the LocationTech
package rename, not a claim that this fork is LocationTech upstream.

\section{Computational geometry issues}

\subsection{Precision model}
JTS lets the client state how many bits of precision are assumed in
input coordinates and maintained in computed coordinates. Two basic
types are supported: Fixed (coordinates on a uniform grid, scale
factor) and Floating (IEEE-754 double or single). Input routines
supplied with JTS round to the precision model. Incompatible
precision models are not reconciled.

In the Fixed model, coordinates fall on a discrete grid. The grid
size is the inverse of the scale factor:
\[
  \mathrm{jtsPt}.x = \mathrm{round}(\mathrm{inputPt}.x \cdot s) / s,
  \qquad
  \mathrm{jtsPt}.y = \mathrm{round}(\mathrm{inputPt}.y \cdot s) / s.
\]
Precise coordinates are represented internally as IEEE-754 doubles
(53 bits of precision).

Floating Double uses the full Java double. Floating Single rounds
computed coordinates to single precision (for example when the
destination is Java2D).

JTS does not implement an exact (rational) precision model.

\subsection{Constructed points and dimensional collapse}
Overlay may construct points that are not input vertices.
Intersection coordinates can require twice the input precision; JTS
rounds them to the PrecisionModel. Rounding can produce dimensional
collapse. JTS forms the lower-dimension Geometry that results.

\subsection{Robustness and numerical stability}
Fundamental predicates (orientation, line intersection, point-in-ring)
use robust algorithms. The binary-predicate algorithm is completely
robust. Overlay and buffer are engineered to return correct answers
in the majority of cases; they are not a substitute for exact
arithmetic. Line-intersection inputs are normalized toward the origin
to preserve bits.

\subsection{Performance --- monotone chains}
Intersection detection uses in-memory spatial indexes and monotone
chains: sequences of segments whose direction vectors lie in one
quadrant. Segments inside one chain do not intersect; envelopes of
contiguous subsets are the envelopes of the endpoints, so binary
search applies. This is unchanged for linear Geometry. It is not a
circular noder.

\section{Spatial model (linear, kept)}
The 2003 design decisions stand for linear Geometry: repeated points
are allowed (SFS); LineStrings may be non-simple (SFS); polygon rings
may not self-touch; polygon rings may mutually touch at a single
point.

\subsection{Geometric definitions (linear)}
Every Geometry carries a PrecisionModel and an optional user-data
field. Empty geometries are allowed; a generic empty is an empty
GeometryCollection. The dimension of a heterogeneous
GeometryCollection is the maximum dimension of its elements.

\paragraph{Curve (SFS, 2003).}
The 2003 spec defined SFS Curve as a non-degenerate linear curve: at
least two points, no two consecutive points equal. The only concrete
curve was LineString / LinearRing. That paragraph is still true of
core.

\paragraph{MultiCurve (SFS, 2003).}
Boundary uses the Mod-2 rule: a point is on the boundary iff it is
on the boundary of an odd number of elements. The 2003 figures
showing that this can disagree with point-set topology still apply
to MultiLineString.

\paragraph{LineString / LinearRing / Polygon / MultiPolygon.}
Unchanged. LineStrings may self-intersect. LinearRings have at least
4 points (including the repeated close), are simple, and may be CW
or CCW. Polygon shells and holes are LinearRings; holes may touch
the shell or another hole at a single point only; interiors are
connected.

\subsection{Support classes}
Coordinate, CoordinateSequence, Envelope, IntersectionMatrix,
GeometryFactory, CoordinateFilter, GeometryFilter remain as in 2003.
GeometryFactory on this tree also has
\texttt{createCircularString} / \texttt{createCompoundCurve} /
\texttt{createCurvePolygon} / \texttt{createMultiCurve} /
\texttt{createMultiSurface} stubs that throw unless the factory is a
\texttt{CurveGeometryFactory} (jts-curve).

\section{Spatial model --- curve types on this tree}
On this tree the concrete SFA/SQL-MM curve types live in the
\texttt{jts-curve} module
(\texttt{org.locationtech.jts.geom.curve}):
\texttt{CircularString}, \texttt{CompoundCurve},
\texttt{CurvePolygon}, and the collections \texttt{MultiCurve} and
\texttt{MultiSurface}. Construction is opt-in: the default
\texttt{GeometryFactory} throws on the curve \texttt{create*}
methods. Use \texttt{CurveGeometryFactory}. Constructors do not
linearise.

\subsection{CircularString (jts-curve)}
A CircularString is a LineString whose consecutive triples (start,
mid, end) are circular arcs. OGC SFA requires an odd point count
$\ge 3$. It is not a LineString of those control chords.
\texttt{toLinear(tolerance)} densifies explicitly. The class is not
constructed by the default GeometryFactory.

\subsection{CompoundCurve (jts-curve)}
A CompoundCurve is a LineString whose members are LineString and/or
CircularString pieces, end-to-start connected. Member structure is
preserved on copy and on \texttt{CurveWKTWriter} /
\texttt{CurveWKBWriter}. A clothoid may appear as a non-leading
member when read by \texttt{CurveWKTReader}; core
\texttt{WKTReader} still fails on the \texttt{CLOTHOID} keyword.
Clothoid membership is extras, not a closed-form construction.

\subsection{CurvePolygon (jts-curve)}
A CurvePolygon is a Polygon whose rings may be CircularString,
CompoundCurve, or LinearRing. A full-circle CircularString shell is
a certified disc. Disc area, envelope, and disc-buffer are closed
form on that shape. A CurvePolygon that is not a certified disc is
not promoted to a laser.

\subsection{MultiCurve and MultiSurface (jts-curve)}
MultiCurve extends MultiLineString; MultiSurface extends
MultiPolygon. They are the SQL-MM collections (WKB 11 and 12). A
MultiSurface of one certified disc unwraps to that disc for the D-HF
disc pair. They are not a general overlay noder.

\subsection{What this section does not add}
No public noder. Constructors do not linearise. This fork is not
described as LocationTech upstream. Triangle / PolyhedralSurface /
TIN exist in the same module as additional ISO types; they are not
the curve overlay stack.

\section{Basic geometric algorithms}
The 2003 primitives remain the linear primitives: point-line
orientation, line intersection test and computation, point-in-ring,
ring orientation. They operate on coordinates and segments. They do
not node circular arcs. Arc/circle predicates used by certified-disc
paths live in \texttt{jts-curve} (package-private helpers) and are
not a public noder.

\section{Topological computation and overlay}
The 2003 spec described topology graphs, labels, the Relate
algorithm, and the overlay algorithm that built a noded graph and
extracted result edges. That is the historical linear overlay.
Production linear overlay on this tree is OverlayNG.

The curve overlay type is \texttt{OverlayNGCurve}. It lives in
\texttt{jts-curve}
(\texttt{org.locationtech.jts.operation.overlayng.curve}). There is
no public noder. Closed-form answers are limited to certified discs.
Anything else densifies and delegates to core OverlayNG.
\texttt{Geometry.intersection} / \texttt{union} / \texttt{difference}
/ \texttt{symDifference} on a curve type are not a general circular
overlay.

Relate / DE-9IM on a certified disc vs Point, LineString, hole-free
Polygon, or a second disc has closed-form cases in the curve module.
Half-disc and general CompoundCurve relate that have no kit fall
through to linearise. That fall-through is not a general arc-aware
relate.

\section{Binary predicates}
Equals, Disjoint, Intersects, Touches, Crosses, Within, Contains,
Overlaps keep their 2003 / SFS meanings on linear Geometry. Two
areas never Crosses. The 2003 contains fix (1.4.1) stands. On curve
operands, a predicate is exact only when a certified path answers;
otherwise it sees chords or an explicit \texttt{toLinear} copy.

\section{Spatial analysis methods}
Constructive methods (Buffer, ConvexHull) and set-theoretic methods
(Intersection, Union, Difference, SymDifference) keep their 2003
point-set definitions. Representation of computed geometries is
still constrained by the precision model.

\paragraph{Buffer (linear).}
Unchanged: Minkowski sum with a disc, approximated by quadrant
segments, end-cap styles as in the Developer's Guide.

\paragraph{Buffer (curve, jts-curve).}
Certified disc $\rightarrow$ disc of radius $r+d$ or empty. Open-arc
buffer is BufferOp on chords --- not a closed form.

\paragraph{ConvexHull (linear).}
Unchanged: smallest convex polygon, or a lower-dimension Geometry if
fewer than 3 hull points; no three consecutive hull points collinear.

\paragraph{ConvexHull (curve, jts-curve).}
A certified disc's hull is the disc (radius-5 disc: area $25\pi$).
A single-arc hull is the slice bounded by the arc and its chord (a
CurvePolygon). The compound-curve hull of circular-plus-straight
members (H-CC), including a stadium, keeps the exposed arcs as
CircularString members of a CurvePolygon. The H-CC witness
\texttt{COMPOUNDCURVE (CIRCULARSTRING (0 0, 5 5, 10 0), (10 0, 10 10))}
has hull area $50 + 12.5\arccos(0.6) \approx 61.59119$ and a
CircularString shell. Any other curve uses the linear hull of
\texttt{toLinear} --- a fallback, not a fail.

\paragraph{Inscribed and empty circles (jts-curve).}
\texttt{MaximumInscribedCircle} on a certified disc and on a stadium
takes a closed form. \texttt{LargestEmptyCircle} on legal curve
obstacles is present. These are constructions, not a general
circular noder.

\fighole{TS-2 disc hull}
\fighole{TS-3 single-arc hull}
\fighole{TS-6 compound hull}

\section{Other methods}
Boundary, isClosed, isSimple, isValid keep the 2003 tables for
linear types. JTS does not validate on construct (efficiency).
IsValidOp reports the nature and location of a failure. Curve
validation is best-effort structural (odd CircularString count,
member connectivity); it is not a full SQL-MM validator.

\section{Well-Known Text and Well-Known Binary}

\subsection{WKT --- linear (kept)}
SFS \S3.2.5. JTS follows the \texttt{MULTIPOINT (x y, x y)} example
form, not the parenthesized-pairs grammar. JTS extends WKT with
\texttt{LINEARRING}. \texttt{WKTReader} is parameterized by a
GeometryFactory (precision and SRID). \texttt{WKTWriter} rounds to
the precision model.

\subsection{WKT --- curve (jts-curve)}
Core \texttt{WKTReader} still rejects \texttt{CIRCULARSTRING} and
the other SQL-MM keywords. Read them with \texttt{CurveWKTReader}.
Write structured members with \texttt{CurveWKTWriter}. Core
\texttt{WKTWriter} already emits the subclass keyword via
\texttt{getGeometryType()} for a flat CircularString; composite
types need the curve writer to keep member tags.

Added tagged text (SQL-MM): \texttt{CIRCULARSTRING},
\texttt{COMPOUNDCURVE}, \texttt{CURVEPOLYGON}, \texttt{MULTICURVE},
\texttt{MULTISURFACE}. \texttt{CLOTHOID} is recognised only by
\texttt{CurveWKTReader}, and only as a non-leading
\texttt{COMPOUNDCURVE} member. That path is extras, not a laser.
Core \texttt{WKTReader} still fails on those keywords.

\subsection{WKB}
The 2003 spec did not list WKB codes 8--12. On this tree:

\begin{center}
\begin{tabular}{clp{4.4cm}p{5.4cm}}
  \toprule
  \textbf{Code} & \textbf{Type} & \textbf{Read (core WKBReader)} &
    \textbf{Write} \\
  \midrule
  1--7 & Point \ldots GeometryCollection &
    Yes, default factory & WKBWriter \\
  8 & CircularString &
    Yes; \texttt{factory.createCircularString} &
    CurveWKBWriter. Bare WKBWriter emits 2 \\
  9 & CompoundCurve &
    Yes; \texttt{factory.createCompoundCurve} &
    CurveWKBWriter. Bare WKBWriter emits 2 \\
  10 & CurvePolygon &
    Yes; \texttt{factory.createCurvePolygon} &
    CurveWKBWriter. Bare WKBWriter emits 3 \\
  11 & MultiCurve &
    Yes; \texttt{factory.createMultiCurve} &
    CurveWKBWriter. Bare WKBWriter emits 5 \\
  12 & MultiSurface &
    Yes; \texttt{factory.createMultiSurface} &
    CurveWKBWriter. Bare WKBWriter emits 6 \\
  \bottomrule
\end{tabular}
\end{center}

WKB type codes 8--12 are first-class in core \texttt{WKBReader} and
\texttt{WKBConstants}. Construction still requires a curve-capable
factory; otherwise the reader throws. \texttt{CurveWKBWriter} emits
codes 8--12. A bare \texttt{WKBWriter} still walks the parent
\texttt{instanceof} ladder and writes a CircularString as LineString
(code 2) and a CurvePolygon as Polygon (code 3). Default and export
WKB paths route through \texttt{CurveWKBWriter}.

Default GeometryFactory throws on the \texttt{create*} curve
methods, so \texttt{new WKBReader().read(type8)} fails with a
factory-required ParseException.
\texttt{new WKBReader(new CurveGeometryFactory())} or
\texttt{CurveWKBReader} succeeds.

\fighole{TS-1 WKB 8--12}

\section{Discrete Hausdorff distance}
The 2003 spec did not define Hausdorff. On this tree the public
class is
\texttt{org.locationtech.jts.algorithm.distance.DiscreteHausdorffDistance}.

Public \texttt{DiscreteHausdorffDistance} (commit \texttt{0ca71b})
has a closed form for two pairs only:
\begin{enumerate}
  \item A single-arc CircularString toward a single-segment
        LineString, apex $\sqrt{949}/6 - 7/6 = 3.967640600249787$ for
        \texttt{CIRCULARSTRING (0 0, 2 3, 10 0)} vs
        \texttt{LINESTRING (0 0, 10 0)}.
  \item Two circular discs (10.0). A
        MultiSurface unwrap of one disc is the same disc pair, not a
        third.
\end{enumerate}
The exact path skips densify; densify is not the laser. For every
other pair the public API still sees vertices / chords. Tests keep
\texttt{fail()} on the general case.

\texttt{DensifyFrac} still interpolates chords on the fallback path.
Calling \texttt{distance(..., densifyFrac)} on a certified pair does
not replace the closed form with densify; the exact path skips
densify. Do not document densify as the laser.

Fr\'echet (\texttt{DiscreteFrechetDistance}) remains open as a
general curve metric. This specification does not claim a Fr\'echet
laser.

\fighole{TS-4 APEX}
\fighole{TS-5 two discs}

\section{Licence}
JTS is dual-licensed under the Eclipse Public License 2.0 and the
Eclipse Distribution License 1.0. See \texttt{LICENSE\_EPLv2.txt},
\texttt{LICENSE\_EDLv1.txt}, and \texttt{LICENSES.md}. The 2005
TestBuilder guide's ACME licence appendix is historical and is not
the project licence.

\section{References}
\begin{itemize}
  \item [SFS] OpenGIS Simple Features Specification For SQL
        Revision 1.1. Open GIS Consortium, Inc. OpenGIS project
        Document 99-049.
  \item [SFA] OpenGIS Implementation Standard for Geographic
        information --- Simple feature access --- Part 1: Common
        architecture, 1.2.1 / ISO 19125.
  \item [Ava97] F.\ Avnaim, J-D.\ Boissonnat, O.\ Devillers,
        F.\ Preparata and M.\ Yvinec. Evaluating signs of
        determinants using single-precision arithmetic.
        \emph{Algorithmica} 17:111--132, 1997.
  \item [Bri98] A.\ Brinkmann, K.\ Hinrichs. Implementing exact line
        segment intersection in map overlay. SDH 1998, pp.\ 569--579.
  \item [Sch97] Stefan Schirra. Precision and robustness in geometric
        computations. In \emph{Algorithmic Foundations of Geographic
        Information Systems}, LNCS 1340, Springer, 1997.
\end{itemize}

\vspace{1.2em}
\noindent\textit{End of Technical Specifications. Named figure holes
are empty 4:3 slots (target 1600$\times$1200, full TestBuilder
window). No clothoid DOI is cited.}

\end{document}
